% =============================================================
%  Chapter — Specification
%  Results of the RTEi / American Terawatt follow-up session on
%  the specification structure (23 Sep 2026, with M. Schudel and
%  H. Saad): maintenance strategy, valve hall and converter
%  requirements, control, protection and data recording.
%  Opens with an abbreviations slide.
%  Second part: next steps and open specification items.
% =============================================================
\sectionslide{Specification}{RTEi follow-up session, 23 Sep 2026}

% ---- Abbreviations ------------------------------------------
\begin{frame}{Abbreviations\\Terms used in this section}
  \noindent{\footnotesize%
  \renewcommand{\arraystretch}{1.15}%
  \begin{tabularx}{\linewidth}{@{}l >{\raggedright\arraybackslash}X@{}}
    \toprule
    \multicolumn{1}{c}{\textbf{Abbr.}} & \textbf{Meaning} \\
    \midrule
    LTSA  & Long-term service agreement \\
    \atrowfill
    MTSI  & Maintenance and technical support \\
    \atrowfill
    O\&M  & Operation and maintenance \\
    \bottomrule
  \end{tabularx}}
\end{frame}


% ---- Maintenance strategy -----------------------------------
\begin{frame}{Maintenance Strategy\\In-house maintenance vs.\ LTSA}

  \begin{itemize}
    \setlength{\itemsep}{0.4em}
    \item RTE maintains in-house from the beginning, supervised by the EPC / MTSI contractor, to build knowledge early.
    \item An MTSI contract (about EUR~3--3.5\,M) is too small relative to the EPC to carry the availability risk; an LTSA limits how far availability can be pushed onto the contractor.
    \item Warranty (3--5~years, hard to extend) covers repair cost; availability is handled by separate liability clauses.
    \item An OEM hotline / deep-expertise contract remains necessary in either case.
  \end{itemize}
\end{frame}


% ---- Maintenance effort -------------------------------------
\begin{frame}{Maintenance Effort\\Staffing and recurring tasks per converter}
  \noindent{\color{atGray}\small RTE reference values for one converter station}
  \vspace{0.6em}

  {\footnotesize
  \renewcommand{\arraystretch}{1.25}
  \begin{tabularx}{\linewidth}{@{}l >{\raggedright\arraybackslash}X >{\raggedright\arraybackslash}X@{}}
    \toprule
    \multicolumn{1}{c}{\textbf{Item}} & \textbf{Scope} & \textbf{Effort} \\
    \midrule
    In-house staff    & Converter operation and maintenance              & About 25 persons \\
    \atrowfill
    Subcontractors    & Auxiliaries (HVAC, diesel generators)            & 3--5 persons \\
    Monthly           & Visual inspection                                & About half a day, on-duty staff \\
    \atrowfill
    Quarterly         & Backups, filter changes                          & About half a day, 2 persons \\
    \atrowhighlight
    Planned outage    & Power module exchange --- critical path, limited by simultaneous lifts & Overnight work to reduce downtime \\
    \bottomrule
  \end{tabularx}}
\end{frame}


% ---- Valve testing ------------------------------------------
\begin{frame}{Valve Testing\\IEC 62501 type tests and site elevation}
  \begin{itemize}
    \setlength{\itemsep}{0.4em}
    \item Dielectric tests: AC, DC and impulse on a valve section with short-circuited submodules; AC--DC between terminals with non-short-circuited submodules.
    \item Wet type test recommended; coolant (deionized water vs.\ glycol) to be negotiated and proven by the contractor.
    \item Existing certificates only for identical designs; repeating tests adds about 6~months to delivery.
    \item Site at about $1500\,\mathrm{m}$: altitude correction of at least $500\,\mathrm{m}$ required --- by test or by calculation is open.
  \end{itemize}
\end{frame}


% ---- Valve hall ---------------------------------------------
\begin{frame}{Valve Hall\\Grounding, hotspots and reactors}
  \begin{itemize}
    \setlength{\itemsep}{0.4em}
    \item \textbf{Grounding} Automatic short-circuiting of the capacitor terminals at tower level with a single switch. Offered by GE, Hitachi and Siemens at a cost premium; discharge still takes 20--30~minutes before entry.
    \item \textbf{Hotspot monitoring} Fixed thermal cameras on the connections towards valves and bushings; robots acceptable if they sweep about once per hour. Tower interiors are too reflective for thermal imaging.
    \item \textbf{Reactor placement} Indoor DC yard under consideration for the desert site. Reactors in the converter hall require thermal studies (heat lake, chimney effect); a separate hall adds bushing-related unavailability, to be quantified in the RAM study.
  \end{itemize}
\end{frame}


% ---- Cooling and replacement --------------------------------
\begin{frame}{Cooling and Module Replacement\\Availability requirements}
  \begin{block}{Key statement}
    Auxiliary failures, cooling in particular, are a leading cause of outages in OEM data --- redundancy must be specified, not left to derating.
  \end{block}
  \vspace{0.3em}
  \begin{itemize}
    \setlength{\itemsep}{0.4em}
    \item Air coolers with $10\,\%$ design margin plus N+1 redundancy; above the temperature limit the redundant cooler is used instead of derating or tripping.
    \item Pump configuration left to the supplier but justified by the RAM study; typical practice is 2~pumps with weekly changeover.
    \item Documentation must demonstrate power module replacement in under 2~hours.
  \end{itemize}
\end{frame}


% ---- Control and protection ---------------------------------
\begin{frame}{Control and Protection\\SCADA split and CMP specification}
  \begin{itemize}
    \setlength{\itemsep}{0.4em}
    \item \textbf{Converter SCADA} Supplied by the OEM. Hitachi is fully integrated in-house; Siemens Energy and GE Vernova use subsuppliers for some components.
    \item \textbf{High-level SCADA} Separate lot: coordination of HVDC and power plants, active power and voltage control, HMI, O\&M.
    \item \textbf{Interface} To be specified in principle; the signal list is defined during design, with wording that allows time for adaptation.
    \item \textbf{Replica system} Likely required.
  \end{itemize}
\end{frame}


% ---- Data recording -----------------------------------------
\begin{frame}{Data Recording\\Fault recorder and trend storage}
  \noindent{\color{atGray}\small Three sampling tiers; a local hard disk must guarantee at least 6~months of storage (RTE moving to 12~months).}
  \vspace{0.6em}

  {\footnotesize
  \renewcommand{\arraystretch}{1.25}
  \begin{tabularx}{\linewidth}{@{}l l >{\raggedright\arraybackslash}X@{}}
    \toprule
    \multicolumn{1}{c}{\textbf{Tier}} & \multicolumn{1}{c}{\textbf{Sampling}} & \textbf{Storage} \\
    \midrule
    Transient     & $50$--$200\,\mu\mathrm{s}$ & Rolling buffer, overwritten unless an event triggers \\
    \atrowfill
    Intermediate  & $5\,\mathrm{ms}$           & Event and dynamic records \\
    \atrowhighlight
    Trend         & $500\,\mathrm{ms}$         & Continuous: temperatures, conductivity and other slow signals \\
    \bottomrule
  \end{tabularx}}

  \vspace{0.6em}
  \noindent{\color{atGray}\small Access via HMI or remote access to the local disk; RTE backs up every 3~months and moved away from remote access for cybersecurity reasons.}
\end{frame}


% =============================================================
%  Next steps
% =============================================================
\intermezzo{Next Steps}{Working packages for AmWatt}


% ---- Open items ---------------------------------------------
\begin{frame}{Open Specification Items\\To be decided before drafting}
  \begin{itemize}
    \setlength{\itemsep}{0.4em}
    \item Definition of the maintenance strategy: in-house vs.\ LTSA
    \item Defenition of the waranty requirements, longer waranty commes with additional costs
    \item Analysis of the NERC cybersecurity requirements to define which communication channels with the HVDC system are possible
    \item Operational data storage definition of the requirements and data retention strategy based on NERC requirements and procedures
    \item Should the automatic sort-circuiting of the capacitor terminals at tower level be requested

  \end{itemize}
\end{frame}

\begin{frame}{Open Specification Items\\To be decided before drafting}
  \begin{itemize}
    \setlength{\itemsep}{0.4em}
    \item DC yard
      \begin{itemize}
        \item indoor vs. outdoor
        \item DC yard shared with converter hall vs. separate indoor DC yard
      \end{itemize}
    \item Type tests
      \begin{itemize}
        \item Definition if type test is required and or corrections can be made by calculation
        \item Definition how to handle the $1500\,\mathrm{m}$ elevation requirement
        \item A type test of the converter could delay the converter delivery by up to 6 months
        \item Definition if wet test procedures are necessary. It was not common in the past, but some operators require those tests now
      \end{itemize}
    \item High-level SCADA interface wording
  \end{itemize}
\end{frame}
